% ---------------------------------------------------------------------------
% Hoja de vida — Manuel Josué Malla Campoverde
% Versión en español adaptada a la convocatoria "Boot Camp Legislativo 2026"
% (Asamblea Nacional del Ecuador). Mantiene el formato del CV maestro en
% inglés (cv/cv.pdf) y reordena el contenido para destacar liderazgo,
% participación ciudadana e impacto público.
%
% Compilar con:  pdflatex cv-bootcamp-legislativo-2026-es.tex
% ---------------------------------------------------------------------------
\documentclass[letterpaper,10pt]{article}

\usepackage[utf8]{inputenc}
\usepackage[T1]{fontenc}
\usepackage{lmodern}        % tipografía vectorial (evita fuentes bitmap)
\usepackage[spanish,es-noshorthands,es-nolists,es-noquoting]{babel}
\usepackage[left=0.7in,right=0.7in,top=0.5in,bottom=0.5in]{geometry}
\usepackage{titlesec}
\usepackage{enumitem}
\usepackage[hidelinks]{hyperref}

\pagestyle{empty}
\setlength{\parindent}{0pt}
\setlength{\tabcolsep}{0pt}

% --- Encabezados de sección: negrita + línea horizontal --------------------
\titleformat{\section}
  {\vspace{2pt}\bfseries\normalsize}
  {}{0em}{}
  [\vspace{-7pt}\rule{\textwidth}{0.9pt}\vspace{-4pt}]
\titlespacing*{\section}{0pt}{6pt}{2pt}

% --- Utilidades ------------------------------------------------------------
% Fila con texto a la izquierda y a la derecha
\newcommand{\fila}[2]{%
  \begin{tabular*}{\textwidth}{@{}l@{\extracolsep{\fill}}r@{}}
    #1 & #2 \\
  \end{tabular*}%
}

% Encabezado de una entrada: cargo/título + fecha, luego organización + lugar
\newcommand{\entrada}[4]{%
  \fila{\textbf{#1}}{#2}\\[1pt]
  \fila{\textit{#3}}{\textit{#4}}\vspace{1pt}%
}

% Lista de viñetas compacta
\newenvironment{vinetas}
  {\begin{itemize}[leftmargin=1.4em,label=$\bullet$,itemsep=0.5pt,topsep=1pt,parsep=0pt,partopsep=0pt]}
  {\end{itemize}\vspace{1pt}}

\begin{document}

% =========================== ENCABEZADO ====================================
\begin{center}
  {\LARGE \textbf{MANUEL JOSUÉ MALLA CAMPOVERDE}}\\[2pt]
  {\small Representante estudiantil $\cdot$ Datos y participación ciudadana}\\[2pt]
  Machala, El Oro, Ecuador $\diamond$ \href{mailto:mmalla1@outlook.com}{mmalla1@outlook.com}
  $\diamond$ +593 97 948 9503 $\diamond$ \href{https://jossuema.tech}{jossuema.tech}
\end{center}
\vspace{-4pt}

% ============================= PERFIL ======================================
\section*{PERFIL}

Estudiante de Ingeniería en Tecnologías de la Información, con interés en el
vínculo entre tecnología, innovación y transformación social. Mi trayectoria
combina el desarrollo de soluciones tecnológicas, la investigación aplicada,
la representación estudiantil ante la Asamblea de Educación Superior de la
Universidad Técnica de Machala y la participación en iniciativas educativas y
comunitarias. Busco fortalecer mis capacidades de participación ciudadana y
adquirir formación legislativa para contribuir al diseño de soluciones que
respondan a las necesidades de los jóvenes y las comunidades del Ecuador.

% =============== LIDERAZGO Y PARTICIPACIÓN CIUDADANA =======================
\section*{LIDERAZGO Y PARTICIPACIÓN CIUDADANA}

\entrada{Representante estudiantil ante la Asamblea de Educación Superior}{noviembre 2025 -- actualidad}
        {Universidad Técnica de Machala (UTMACH)}{Machala, Ecuador}
\begin{vinetas}
  \item Represento a los estudiantes en espacios institucionales de educación
        superior.
  \item Gestiono alianzas con instituciones públicas y privadas para
        respaldar iniciativas educativas estudiantiles.
\end{vinetas}

\entrada{Docente voluntario de fundamentos de programación}{marzo -- abril 2025}
        {Municipio de Huaquillas}{Huaquillas, Ecuador}
\begin{vinetas}
  \item Impartí un curso introductorio de programación en Python a jóvenes
        estudiantes, en alianza con el gobierno local, para fomentar la
        alfabetización digital en la comunidad.
\end{vinetas}

\entrada{Embajador estudiantil de Microsoft Learn}{enero 2024 -- actualidad}
        {Microsoft}{Machala, Ecuador}
\begin{vinetas}
  \item Promuevo la comunidad tecnológica y el aprendizaje entre pares
        mediante talleres de inteligencia artificial y ciencia de datos.
\end{vinetas}

\entrada{Voluntario organizador --- conferencia PhawAI y TaReCDa 2025}{junio -- noviembre 2025}
        {Comunidad de inteligencia artificial y ciencia de datos}{Arequipa, Perú}
\begin{vinetas}
  \item Apoyé la organización de la conferencia de inteligencia artificial y
        ciencia de datos PhawAI + TaReCDa 2025.
  \item Colaboré con el equipo organizador del reto de Somos NLP, comunidad
        hispanohablante de procesamiento de lenguaje natural (mayo 2025).
\end{vinetas}

\entrada{Miembro del equipo de intercambios de voluntariado internacional}{febrero -- junio 2024}
        {AIESEC en Ecuador}{Machala, Ecuador}
\begin{vinetas}
  \item Apoyé el reclutamiento y la coordinación de voluntarios para
        proyectos internacionales vinculados a los Objetivos de Desarrollo
        Sostenible (ODS).
\end{vinetas}

% =========================== EDUCACIÓN =====================================
\section*{EDUCACIÓN}

\entrada{Ingeniería en Tecnologías de la Información}{2022 -- 2027}
        {Universidad Técnica de Machala}{Machala, Ecuador}
Asignaturas relevantes: probabilidad y estadística, análisis matemático,
métodos numéricos, bases de datos, ingeniería de software y seguridad de la
información.
\vspace{2pt}

\entrada{Especialista en Decisiones Basadas en Datos}{agosto 2025 -- marzo 2026}
        {Escuela Superior Politécnica del Litoral (ESPOL)}{Guayaquil, Ecuador}
Programa cursado con beca. Asignaturas: SQL y optimización de consultas,
Python para analítica de datos, inteligencia de negocios, visualización de
datos e IA generativa aplicada a la analítica.
\vspace{2pt}

\entrada{Programa de verano en matemática aplicada}{semestre 2024-II}
        {Escuela de Matemática Aplicada de la FGV (FGV EMAp)}{Río de Janeiro, Brasil}
Asignaturas: aprendizaje automático aplicado a sistemas de recomendación,
métodos de optimización en Julia, redes complejas aplicadas a la epidemiología
y regresión bayesiana aplicada.

% ====================== EXPERIENCIA PROFESIONAL ============================
\section*{EXPERIENCIA PROFESIONAL}

\entrada{Ingeniero de aprendizaje automático (medio tiempo)}{noviembre 2025 -- actualidad}
        {LarvIA}{Machala, Ecuador}
\begin{vinetas}
  \item Desarrollo, entrenamiento y despliegue de modelos de visión por
        computadora en equipos que procesan la información directamente en
        campo, sin conexión permanente a internet.
\end{vinetas}

\entrada{Pasante de investigación en aprendizaje automático y ciencia de datos}{enero -- febrero 2026}
        {CNPEM --- Laboratorio Nacional de Nanotecnología del Brasil}{Campinas, Brasil}
\begin{vinetas}
  \item Análisis de datos de sensores electroquímicos y desarrollo de modelos
        para detectar patrones en sus mediciones, en un laboratorio nacional
        de investigación.
\end{vinetas}

\entrada{Pasante de ingeniería de aprendizaje automático}{mayo -- octubre 2025}
        {LarvIA}{Machala, Ecuador}
\begin{vinetas}
  \item Diseño y despliegue de sistemas de análisis automatizado de imágenes
        para la optimización del cultivo de camarón en el Ecuador y en
        Latinoamérica, sector clave para el empleo y las exportaciones del
        país.
  \item Aplicación de técnicas de fotogrametría para obtener medidas reales a
        partir de fotografías tomadas con teléfonos móviles.
\end{vinetas}

\entrada{Desarrollador full-stack junior}{febrero -- mayo 2025}
        {Ejercicio profesional independiente}{Machala, Ecuador}
\begin{vinetas}
  \item Desarrollo de un panel de administración empresarial y de un sistema
        de gestión de información, junto con un asistente conversacional para
        el agendamiento de reuniones a través de WhatsApp.
\end{vinetas}

\entrada{Pasante de desarrollo de software}{octubre 2024 -- febrero 2025}
        {Universidad Técnica de Machala --- Departamento de TI}{Machala, Ecuador}
\begin{vinetas}
  \item Diseño de la arquitectura del Sistema Inteligente de Control de
        Acceso y Seguridad (SISCA) para las aulas de las facultades del campus
        principal de la universidad pública.
  \item Diseño de la aplicación móvil utilizada por docentes y personal
        administrativo para el acceso a las aulas.
\end{vinetas}

% ======================== PROYECTOS DE IMPACTO =============================
\section*{PROYECTOS DE IMPACTO}

\begin{vinetas}
  \item \textbf{Análisis de conglomerados geoespaciales de homicidios en el
        Ecuador: comparación de los métodos KMeans, DBSCAN y HDBSCAN} (2025)\\
        Investigación desarrollada a partir de datos gubernamentales de
        homicidios con el fin de identificar conglomerados territoriales y
        zonas donde el fenómeno se concentra con mayor intensidad y sigue
        patrones similares. Es un ejercicio de análisis de datos públicos con
        potencial aplicación al estudio de problemas de seguridad ciudadana y
        a la focalización territorial de políticas.

  \item \textbf{Buzón Inteligente UTMACH: plataforma de participación
        universitaria}\\
        Plataforma web para la recepción de quejas, sugerencias y propuestas
        de la comunidad universitaria, con un tablero analítico para su
        seguimiento. El proyecto busca acercar los mecanismos de participación
        a los estudiantes y sistematizar sus opiniones para el análisis
        institucional.

  \item \textbf{MangroveNet: mapeo y monitoreo de ecosistemas de manglar}
        (2024)\\
        Modelo de aprendizaje profundo para la identificación de manglares a
        partir de imágenes satelitales Landsat y Sentinel-2. Permite
        monitorear a bajo costo un ecosistema costero estratégico para el
        Ecuador, con aplicaciones en conservación y control ambiental.
\end{vinetas}

% =================== RECONOCIMIENTOS Y DISTINCIONES ========================
\section*{RECONOCIMIENTOS Y DISTINCIONES}

\begin{vinetas}
  \item Premio a la mejor presentación en el hackatón de proyectos --- EPIC
        VI, Escuela de Programación para la Investigación Científica (Quito,
        Ecuador, 2026).
  \item Primer lugar --- UTMACH Challenge, Universidad Técnica de Machala
        (Machala, Ecuador, 2024).
  \item Primer lugar --- Fall AI Project Competition, Microsoft Learn Student
        Ambassadors (2024).
  \item Resumen extendido seleccionado y premio a la mejor presentación en la
        sesión de estudiantes de pregrado --- LatinX in CV, congreso CVPR
        (Seattle, Estados Unidos, 2024).
  \item Propuesta de investigación seleccionada y beca de participación ---
        TaReCDa/ReWARDS (Machala, Ecuador, 2023).
  \item Becas competitivas de formación y participación: programa Especialista
        en Decisiones Basadas en Datos, ESPOL (Guayaquil, Ecuador); beca de
        viaje y alojamiento, FGV EMAp (Río de Janeiro, Brasil, 2025); beca de
        participación, KHIPU 2025, Encuentro Latinoamericano de Inteligencia
        Artificial (Santiago, Chile); beca de viaje e inscripción, Computer
        Vision Foundation (Seattle, Estados Unidos, 2024).
\end{vinetas}

% ============================ HABILIDADES ==================================
\section*{HABILIDADES}

\begin{tabular*}{\textwidth}{@{}p{0.22\textwidth}@{}p{0.78\textwidth}@{}}
  \textbf{Competencias} & Representación estudiantil, facilitación de talleres,
                          trabajo en equipos multiculturales y comunicación de
                          resultados técnicos a públicos no especializados. \\[3pt]
  \textbf{Datos} & Análisis de datos públicos y gubernamentales, estadística
                   aplicada, análisis geoespacial y visualización de datos
                   (Power BI). \\[3pt]
  \textbf{Tecnología} & Python, SQL, aprendizaje automático (\textit{machine
                        learning}) y visión por computadora. \\[3pt]
  \textbf{Idiomas} & Español (nativo), inglés (C1), portugués (B1). \\
\end{tabular*}

\end{document}
